\documentclass[11pt,a4paper]{ltjsarticle}
\usepackage{luatexja}
\usepackage{luatexja-fontspec}
\usepackage{amsmath,amssymb}
\usepackage{graphicx}
\usepackage{longtable,booktabs,array}
\usepackage{float}
\usepackage[margin=25mm]{geometry}
\usepackage{caption}
\usepackage{hyperref}
\hypersetup{colorlinks=true,linkcolor=blue,urlcolor=blue,citecolor=blue}
\setkeys{Gin}{width=0.82\linewidth,keepaspectratio}
\usepackage{amsthm}
\newtheorem{theorem}{定理}
\newtheorem{proposition}{命題}
\newtheorem{lemma}{補題}
\newtheorem{corollary}{系}
\newtheorem{definition}{定義}
\newtheorem{remark}{注意}
\newtheorem{observation}{観察}
\usepackage{calc}
\newcounter{none}
\providecommand{\real}[1]{#1}
\providecommand{\pandocbounded}[1]{#1}
\providecommand{\tightlist}{\setlength{\itemsep}{0pt}\setlength{\parskip}{0pt}}
\title{共役複素ノルムと平方量読出しの接続について}
\author{木原 範昭 (Noriaki Kihara), WF System Co., Ltd.~--- DOI:
10.5281/zenodo.21127200}
\date{v0.5 --- 2026年7月}
\begin{document}
\maketitle
\emph{位相表示から不可視軸を含む多次元二乗和表示への拡張観察 ------
初等代数と幾何的解釈に限った非物理的整理}

\section{0.
前提と立ち位置（本稿の読み方）}\label{ux524dux63d0ux3068ux7acbux3061ux4f4dux7f6eux672cux7a3fux306eux8aadux307fux65b9}

本稿は、\textbf{共役な複素数とその幾何学的解釈のみ}から出発する。物理学（時空、特殊・一般相対論、量子論、計量の署名）とは\textbf{接続しない立場}をとる。したがって、次を明確に宣言する。

\begin{enumerate}
\def\labelenumi{\arabic{enumi}.}
\tightlist
\item
  以下に現れる座標記号 \(x,\ y,\ z_1,\ z_2,\ z_3,\dots\)
  は、すべて\textbf{幾何学的な座標ラベル}である。\textbf{物理的な空間次元・時間次元とは同定しない。}
\item
  特に、本稿では時間を表す記号 \(t\)
  を\textbf{用いない}。虚数値をとる座標を導入する場合も、それは「時間」ではなく、単に\textbf{虚数値をとる座標ラベル}
  \(z_h\) である。
\item
  二乗差的な符号（負号）が式に現れても、それは\textbf{計量の署名を主張するものではない}。負号は、ある座標を虚数として置いたときの
  \(i^2=-1\) に由来する\textbf{代数的帰結}にすぎない。
\item
  すべての座標軸は\textbf{対等}に扱う。ここでいう対等性とは、\textbf{最初から特定の物理的役割や計量署名を割り当てない}という意味である。後の節で、代数的表示として実数値成分と虚数値成分に分けること（§5）は、この対等性に反しない。
\item
  本稿は新しい物理法則・新しい数学定理を主張しない。既知の共役複素ノルムと平方和を、先行する平方量読出し系列の中に置いて読み替える\textbf{観察論文}である。
\end{enumerate}

この宣言は、v0.1
における誤り（特定の軸に手で負号を与え、半径記号を二重に用いたこと）を除くための、本稿の設計上の土台である。

\begin{center}\rule{0.5\linewidth}{0.5pt}\end{center}

\section{要旨}\label{ux8981ux65e8}

本稿は、平方量読出しによる線形単体化、半波長奇数倍音和による孤立ピーク波、二コピー共通相対位相重ね合わせにおけるコントラスト則、の三篇に続く第四篇である。

先行する平方量読出し論文では、正の領域で \(X_i=x_i^2\)
と読むことにより、平方和制約 \(\sum_i x_i^2=E\)
が平方量空間の線形単体制約 \(\sum_i X_i=E\)
として読めることを整理した。本稿では、この平方量読出しに、共役複素数の最も基本的なノルム

\[
Z\bar Z=x^2+y^2
\]

を接続する。複素極表示
\(Z=\rho e^{i\theta}=\rho(\cos\theta+i\sin\theta)\) に
\(x=\rho\cos\theta,\ y=\rho\sin\theta\) を置けば
\(Z=x+iy\)、\(Z\bar Z=x^2+y^2=\rho^2\)
が成り立つ。これは標準的事実であり、新規性はこの式そのものにない。本稿が観察するのは、この既知のノルムが、\textbf{二つの直交実成分の平方量を一つの実ノルムへ束ねる操作}として、平方量読出しと自然に接続する点である。

次に、これを\textbf{対等な多次元拡張}へ一般化する。実成分を二本ずつ複素数に束ねれば

\[
\sum_j Z_j\bar Z_j=\sum_n x_n^2=R^2
\]

となり、複素平面に表示されない座標軸（本稿でいう不可視軸）も、\textbf{同じ符号
\(+\) を持つ対等な軸} \(z_1,z_2,\dots\)
として加わるだけである。ここに特別な軸はない。

最後に、前半の共役ノルム
\(\sum_j Z_j\bar Z_j\)（Hermitian・正定値）に、虚数値座標の二乗
\((i z_k)^2\) を加えた\textbf{零二乗和形式}

\[
\sum_j Z_j\bar Z_j+\sum_k (i z_k)^2=0
\]

を観察する。ここで共役積 \(X\bar X\) と非共役の二乗 \(X^2\)
は、\textbf{実座標では一致するが虚座標でのみ分岐}し（\(X_h\bar X_h=z_h^2\)
に対し \(X_h^2=-z_h^2\)）、この分岐が負号の唯一の起源である。展開すれば
\(\sum_j Z_j\bar Z_j=\sum_k z_k^2\)、すなわち

\[
\sum_n x_n^2=\sum_k z_k^2
\]

が得られる。ここで現れる負号は \(i^2=-1\)
に由来する代数的帰結であって、手で置いた計量の署名ではない。共役ノルム型
\(Z\bar Z\) と虚座標の非共役二乗 \(X_h^2\)
は同一の演算ではなく、本稿はこの区別を保つ。これらはすべて、共役複素数と二乗和の\textbf{代数的配置}に関する幾何学的観察であり、物理学的な時空・相対論の主張を含まない。

\begin{center}\rule{0.5\linewidth}{0.5pt}\end{center}

\section{1.
系列上の位置づけ}\label{ux7cfbux5217ux4e0aux306eux4f4dux7f6eux3065ux3051}

本稿は、中心投影幾何そのものを扱う論文ではない。中心投影シリーズでは、曲率半径・測地セル・角度面積体積補正を主題とし、意図的に共役複素数表現から距離を取っていた。そのため、本稿の内容を中心投影シリーズへ直接挿入すると論旨の重心がずれる。

本稿が属するのは次の系列である。

\begin{enumerate}
\def\labelenumi{\arabic{enumi}.}
\tightlist
\item
  平方量読出しによる線形単体化と曲率補正候補の次元別整理。
\item
  半波長位相区間における一定振幅奇数倍音和の孤立ピーク波。
\item
  二コピー共通相対位相重ね合わせにおける波形不変性とコントラスト則。
\item
  本稿、すなわち共役複素ノルムと平方量読出しの接続、およびその対等な多次元拡張。
\end{enumerate}

本稿は、これらの後に、位相表示・共役複素数・平方量読出しを結ぶ最小の接点として
\(Z\bar Z=x^2+y^2\) を扱い、それを対等な多次元二乗和へ一般化する。

\begin{center}\rule{0.5\linewidth}{0.5pt}\end{center}

\section{2.
複素極表示と二乗和}\label{ux8907ux7d20ux6975ux8868ux793aux3068ux4e8cux4e57ux548c}

\subsection{2.1 複素極表示}\label{ux8907ux7d20ux6975ux8868ux793a}

複素数 \(Z\) を \(Z(\rho,\theta)=\rho(\cos\theta+i\sin\theta)\)
と書く。\(\rho\ge0\) は複素平面上の半径、\(\theta\)
は位相角である。\(x=\rho\cos\theta,\ y=\rho\sin\theta\) と置けば
\(Z=x+iy\) である。共役複素数は \(\bar Z=x-iy\) であるから、

\[
Z\bar Z=(x+iy)(x-iy)=x^2+y^2 .
\]

極表示からも
\(Z\bar Z=\rho e^{i\theta}\rho e^{-i\theta}=\rho^2\)。したがって

\[
Z\bar Z=x^2+y^2=\rho^2 .
\]

\subsection{2.2
平方量読出しとしての解釈}\label{ux5e73ux65b9ux91cfux8aadux51faux3057ux3068ux3057ux3066ux306eux89e3ux91c8}

先行する平方量読出し論文では、正の領域で \(X_i=x_i^2\)
と読み、平方和制約 \(\sum_i x_i^2=E\) を線形制約 \(\sum_i X_i=E\)
として整理した。本稿の \(Z\bar Z=x^2+y^2\)
は、この平方量読出しの\textbf{二成分版}であり、二つの平方量
\(x^2,\ y^2\) が共役複素積 \(Z\bar Z\)
として一つの実ノルムに束ねられている。

\subsection{\texorpdfstring{2.3 二種類の \(i\)
の役割（本稿での区別）}{2.3 二種類の i の役割（本稿での区別）}}\label{ux4e8cux7a2eux985eux306e-i-ux306eux5f79ux5272ux672cux7a3fux3067ux306eux533aux5225}

本稿では、虚数単位 \(i\)
の二つの異なる役割を明確に区別する。この区別を怠ると、後述の多次元拡張で誤りが生じる。

\begin{itemize}
\tightlist
\item
  \textbf{役割①：束ね（bundling）}。\(Z=x+iy\) で \(x,y\)
  がともに実数のとき、\(Z\bar Z=x^2+y^2>0\)。ここで \(i\)
  は、二本の実座標を一つの複素座標へまとめる記法上の道具にすぎない。ノルムは常に非負。
\item
  \textbf{役割②：虚数値座標指定（imaginary-valued coordinate
  label）}。座標ラベルそのものを虚数値として指定するとき、\(X_h=i\,z_h\)（\(z_h\)
  は実ラベル）と置く。このとき \(X_h^2=(i z_h)^2=-z_h^2\)
  となり、\textbf{負の二乗寄与}が生じる。
\end{itemize}

役割①は \(x^2+y^2\)（正定値）を与え、役割②は
\(-z_h^2\)（負寄与）を与える。両者は別物であり、地続きに扱ってはならない。より正確には、役割①は\textbf{共役を含む積}
\(Z\bar Z\)（Hermitian・正定値）に現れ、役割②は\textbf{共役を含まない二乗}
\(X_h^2\)（複素双線形形式）に現れる。この\textbf{共役の有無}の違いが、§5
で決定的になる。

\subsection{2.4
新規性の位置}\label{ux65b0ux898fux6027ux306eux4f4dux7f6e}

\(Z\bar Z=x^2+y^2\)
は標準的事実である。本稿はこの式を新定理として提示しない。目的は、既知の複素ノルムを平方量読出し系列に置き、

\[
\text{平方量読出し}\quad\leftrightarrow\quad\text{共役複素ノルム}
\]

という接点を明示し、それを対等な多次元二乗和へ一般化することである。

\begin{center}\rule{0.5\linewidth}{0.5pt}\end{center}

\section{3.
位相表示としての読み}\label{ux4f4dux76f8ux8868ux793aux3068ux3057ux3066ux306eux8aadux307f}

複素極表示では \(\theta\)
は平面上の角度であり、同時に波の位相角として読むこともできる。すなわち

\[
\operatorname{Re}Z=\rho\cos\theta,\qquad \operatorname{Im}Z=\rho\sin\theta
\]

であり、\(Z=\rho e^{i\theta}\) は位相 \(\theta\)
を持つ状態を複素平面上の直交二成分へ展開した表示である。

本稿では、この段階で扱っているのは
\((\rho,\theta)\longmapsto(x,y)\)、すなわち位相角を含む極表示から二つの直交成分への写像である。ここに時間や進行方向の座標は導入しない（§0
の宣言のとおり、時間記号 \(t\) は用いない）。

\begin{center}\rule{0.5\linewidth}{0.5pt}\end{center}

\section{4.
対等な多次元二乗和への拡張}\label{ux5bfeux7b49ux306aux591aux6b21ux5143ux4e8cux4e57ux548cux3078ux306eux62e1ux5f35}

\subsection{4.1
複数の複素数による束ね}\label{ux8907ux6570ux306eux8907ux7d20ux6570ux306bux3088ux308bux675fux306d}

実座標を \(x_1,x_2,\dots,x_{2m}\) とし、二本ずつ複素数に束ねる。

\[
Z_j=x_{2j-1}+i\,x_{2j},\qquad Z_j\bar Z_j=x_{2j-1}^2+x_{2j}^2 .
\]

これらの和は

\[
\sum_{j=1}^{m} Z_j\bar Z_j=\sum_{n=1}^{2m} x_n^2 .
\]

奇数本の場合は、束ねられない実座標 \(x_{2m+1}\)
を単独で加える。いずれの場合も、

\[
\sum_n x_n^2=R^2
\]

は、全座標を対等に含む正定値の二乗和である。

\subsection{4.2
不可視軸は特別な軸ではない}\label{ux4e0dux53efux8996ux8ef8ux306fux7279ux5225ux306aux8ef8ux3067ux306fux306aux3044}

ある座標が、いま表示している複素平面（たとえば \(Z_1\)
の平面）に現れないとき、本稿ではそれを\textbf{不可視軸}と呼ぶ。しかしこれは、表示上その平面に載っていないという意味にすぎず、\textbf{符号も役割も他の軸と同じ}である。不可視軸
\(z_1,z_2,\dots\) を加えることは、

\[
\sum_n x_n^2 + z_1^2 + z_2^2 + \cdots = R^2
\]

のように、左辺に\textbf{同符号 \(+\) の項を加えて \(R\)
が増える}だけである。半径 \(R\)
は最後まで「全座標の二乗和の平方根」ただ一つであり、二重には定義されない。

ここが v0.1
の誤りを除く要点である。特定の軸だけを取り出して反対側に置いたり、負号を与えたりすると、\(\sum_n x_n^2=R^2\)
が崩れ、\(R\)
が二通りの意味を持ってしまう。対等に扱う限り、その破綻は生じない。

なお、ここでいう「不可視」とは、選んだ複素平面表示に現れないという\textbf{表示上の意味に限る}。物理的に観測不能な軸、隠れた物理次元、または未知の自由度を意味しない。

\subsection{4.3
傾き角による表示（球面座標）}\label{ux50beux304dux89d2ux306bux3088ux308bux8868ux793aux7403ux9762ux5ea7ux6a19}

三本の場合、半径 \(R\) と傾き角 \(\phi\)、位相角 \(\theta\) を用いれば

\[
x=R\cos\phi\cos\theta,\qquad y=R\cos\phi\sin\theta,\qquad z_1=R\sin\phi
\]

と書け、直ちに

\[
x^2+y^2+z_1^2
=R^2\cos^2\phi(\cos^2\theta+\sin^2\theta)+R^2\sin^2\phi
=R^2
\]

となる。ここで \(\phi\)
は複素平面から外れる傾き角であり、平面内の追加回転位相ではない。両者は区別しなければならない。平面内位相としての追加角
\(\alpha\) は \(x=\rho\cos(\theta+\alpha),\ y=\rho\sin(\theta+\alpha)\)
を与えるだけで \(z_1\) 成分を生成せず、\(x^2+y^2=\rho^2\) のままである。

\begin{center}\rule{0.5\linewidth}{0.5pt}\end{center}

\section{5.
共役ノルム型平方量と虚数値座標による零二乗和}\label{ux5171ux5f79ux30ceux30ebux30e0ux578bux5e73ux65b9ux91cfux3068ux865aux6570ux5024ux5ea7ux6a19ux306bux3088ux308bux96f6ux4e8cux4e57ux548c}

\subsection{5.1 共役積と二乗の違い ------
実座標では一致、虚座標で分岐}\label{ux5171ux5f79ux7a4dux3068ux4e8cux4e57ux306eux9055ux3044-ux5b9fux5ea7ux6a19ux3067ux306fux4e00ux81f4ux865aux5ea7ux6a19ux3067ux5206ux5c90}

前節までの \(\sum_j Z_j\bar Z_j=\sum_n x_n^2=R^2\)
は、\textbf{共役を含む正定値ノルム}（Hermitian 形式
\(\sum_n X_n\bar X_n=\sum_n|X_n|^2\)）であり、零になるのは全成分が零のときに限る。

ここで、共役積 \(X\bar X\) と、共役を含まない二乗 \(X^2\)
の関係を明確にする。\textbf{実数値座標では両者は一致する}：\(x\)
が実なら
\(x\bar x=x^2\)。\textbf{虚数値座標でのみ分岐する}：\(X_h=i\,z_h\)（\(z_h\)
は実）に対して、

\[
X_h\bar X_h=(i z_h)(\overline{i z_h})=(i z_h)(-i z_h)=z_h^2\ (>0),
\qquad
X_h^2=(i z_h)^2=-z_h^2\ (<0).
\]

すなわち、\textbf{負の二乗寄与は、虚数値座標に対して共役を含まない二乗をとったときに''だけ''生じる}。これが以降に現れる負号の唯一の起源であり、手で置いた計量署名ではない。したがって本稿は、共役ノルム型
\(Z\bar Z\) と、虚座標の非共役二乗 \(X_h^2\)
を、\textbf{同一の演算として混同しない}。

\subsection{5.2
零二乗和形式（共役ノルム背骨＋虚座標二乗）}\label{ux96f6ux4e8cux4e57ux548cux5f62ux5f0fux5171ux5f79ux30ceux30ebux30e0ux80ccux9aa8ux865aux5ea7ux6a19ux4e8cux4e57}

共役ノルムで得た正定値平方量 \(\sum_j Z_j\bar Z_j\) に、虚数値座標の二乗
\((i z_k)^2\) を加えた\textbf{零二乗和形式}

\[
\sum_j Z_j\bar Z_j+\sum_k (i z_k)^2=0,
\qquad Z_j=x_{2j-1}+i x_{2j},\ \ z_k\in\mathbb{R}
\]

を観察する。ここで可視部は共役ノルム（Hermitian・正定値）、不可視部は虚数値座標の二乗（非共役）であり、両者は異なる演算である。左辺を展開すると
\((i z_k)^2=-z_k^2\) により

\[
\sum_j Z_j\bar Z_j-\sum_k z_k^2=0
\quad\Longrightarrow\quad
\sum_j Z_j\bar Z_j=\sum_k z_k^2 .
\]

さらに \(\sum_j Z_j\bar Z_j=\sum_n x_n^2\) であるから、

\[
\sum_n x_n^2=\sum_k z_k^2 .
\]

なお、可視座標は実数なので §5.1 により \(x_n\bar x_n=x_n^2\)
である。よってこの形式は、全座標を一括して \(\sum_n X_n^2=0\)
と書いた非共役二次形式と数値的には一致する。ただし本稿は、可視部を共役ノルム、不可視部を虚座標二乗として\textbf{分けて}書く上の形式を主とする。\(Z\bar Z\)
型（共役）と \(X^2\) 型（非共役）を混同しないためである。

\subsection{5.3
非自明解は非実成分を要求する}\label{ux975eux81eaux660eux89e3ux306fux975eux5b9fux6210ux5206ux3092ux8981ux6c42ux3059ux308b}

もし全座標が実数なら、\(\sum_j Z_j\bar Z_j=\sum_n x_n^2=0\)
は全零解しか持たない（正定値ノルム）。したがって、零二乗和が\textbf{非自明であるためには、少なくとも一つの非実成分が必要}である。本稿ではその最小整理として、非実成分を純虚数値座標
\(i z_k\) に限って扱う（一般には \(a_n+i b_n\)
の形もありうる）。非実成分は追加の仮定ではなく、非自明性の必要条件として現れる。

\subsection{5.4
負号の由来と非同定}\label{ux8ca0ux53f7ux306eux7531ux6765ux3068ux975eux540cux5b9a}

ここで現れる負号は、§5.1 のとおり虚数値座標の二乗 \(i^2=-1\)
から生じる代数的帰結であって、特定の軸に手で与えた計量署名ではない。また本稿は、\(z_k\)
を時間座標・曲率半径・投影半径のいずれとも同定しない（§0, §6）。

\subsection{5.5 v0.1
の誤りとの対照}\label{v0.1-ux306eux8aa4ux308aux3068ux306eux5bfeux7167}

v0.1 は、正定値の三次元和 \(x^2+y^2+z^2=r^2\) に対し、右辺へ別記号の二乗
\(t^2\) を手で加えて \(x^2+y^2+z^2=r^2+t^2\)
とした。しかしこれは、対等な次元拡張（左辺に同符号項を足す操作）ではなく、（i）新しい座標の導入と（ii）その負号での配置という\textbf{二つの独立操作を混ぜたもの}であった。結果として
\(\sum_n x_n^2=R^2\) が崩れ、\(r\) が二通りの意味を持った。

本稿の形式では、負号は虚座標二乗（\(i^2\)）から\textbf{導かれ}、可視部の共役ノルム
\(Z\bar Z\)（正定値）と不可視部の虚座標二乗 \(X_h^2\)
を明確に分離し、\(R\) を二重に用いないため、この破綻は生じない。

\begin{center}\rule{0.5\linewidth}{0.5pt}\end{center}

\section{6.
本稿が現れさせない読み}\label{ux672cux7a3fux304cux73feux308cux3055ux305bux306aux3044ux8aadux307f}

§0 の宣言を、具体式に即して繰り返す。

\begin{itemize}
\tightlist
\item
  本稿は、時間座標を導入しない。\(\sum_a x_a^2=\sum_k z_k^2\) の右辺
  \(z_k\) は虚数値座標ラベルであり、時間ではない。
\item
  本稿は、計量の署名 \((+,+,\dots,-)\) を主張しない。負の二乗寄与は
  \(i^2=-1\) に由来する代数であって、時空計量ではない。
\item
  本稿は、光円錐・Minkowski
  空間・特殊相対論・一般相対論のいずれも導出しない。\(\sum_a x_a^2=\sum_k z_k^2\)
  が形式的にそれらの二次形式と一致する場合があっても、本稿はその同定を行わない。
\item
  本稿は、\(z_k\) や \(R\)
  を曲率半径・投影半径・内部半径と同一視しない。
\item
  本稿は、座標ラベル \(x,y,z_1,z_2,\dots\) を物理的次元と同定しない。
\item
  本稿でいう「虚数値座標」および「不可視軸」は、座標ラベルの代数的・表示上の指定であって、観測不能な物理自由度や隠れた物理次元を意味しない。
\end{itemize}

これらは、代数的観察を物理主張へ飛躍させないための境界条件である。

\begin{center}\rule{0.5\linewidth}{0.5pt}\end{center}

\section{7.
先行三篇との関係}\label{ux5148ux884cux4e09ux7bc7ux3068ux306eux95a2ux4fc2}

\subsection{7.1
平方量読出し論文との関係}\label{ux5e73ux65b9ux91cfux8aadux51faux3057ux8ad6ux6587ux3068ux306eux95a2ux4fc2}

平方量読出し論文が \((x_i)\longmapsto(x_i^2)\)
という個別平方量読出しを扱ったのに対し、本稿は

\[
(x,y)\longmapsto Z=x+iy\longmapsto Z\bar Z=x^2+y^2
\]

という、共役複素ノルムによる二成分平方和読出しと、その対等な多次元和
\(\sum_j Z_j\bar Z_j=\sum_n x_n^2\) を扱う。

\subsection{7.2
奇数倍音孤立ピーク論文との関係}\label{ux5947ux6570ux500dux97f3ux5b64ux7acbux30d4ux30fcux30afux8ad6ux6587ux3068ux306eux95a2ux4fc2}

第二篇が「位相区間上にどのような局在形が作れるか」を扱ったのに対し、本稿は「位相角を持つ複素表示が、平方量読出しとどう接続するか」を扱う。具体的な奇数倍音和
\(S_N(\varphi)\) ではなく、位相表示 \(Z=\rho e^{i\theta}\)
そのものを対象とする。

\subsection{7.3
相対位相コントラスト論文との関係}\label{ux76f8ux5bfeux4f4dux76f8ux30b3ux30f3ux30c8ux30e9ux30b9ux30c8ux8ad6ux6587ux3068ux306eux95a2ux4fc2}

第三篇では、二コピーに共通相対位相を与えると
\(\psi_\alpha=2\cos\alpha\,S_N\)、\(I_\alpha=4\cos^2\alpha\,I_N\)
となり、相対位相が波形を変えず二乗強度係数として現れた。本稿では、共役複素積
\(Z\bar Z\) によって位相が消え、ノルム平方 \(\rho^2\)
が残る。両者に共通するのは、位相そのものではなく、\textbf{位相を含む構造を二乗量として読んだときに実数量が現れる}点である。

\begin{center}\rule{0.5\linewidth}{0.5pt}\end{center}

\section{8.
危険な読みと安全な読み}\label{ux5371ux967aux306aux8aadux307fux3068ux5b89ux5168ux306aux8aadux307f}

\subsection{8.1
危険な読み（本稿は主張しない）}\label{ux5371ux967aux306aux8aadux307fux672cux7a3fux306fux4e3bux5f35ux3057ux306aux3044}

\begin{itemize}
\tightlist
\item
  複素数から空間や時空を導出した。
\item
  位相空間と実空間は同一である。
\item
  複素ノルムから特殊相対論・一般相対論を導出した。
\item
  負の二乗寄与は時空計量の署名である。
\item
  座標ラベル \(z_k\) は時間である。
\item
  \(R\) や \(z_k\) は曲率半径である。
\item
  座標ラベルは物理的次元である。
\end{itemize}

\subsection{8.2
安全な読み（本稿で言えること）}\label{ux5b89ux5168ux306aux8aadux307fux672cux7a3fux3067ux8a00ux3048ux308bux3053ux3068}

\begin{enumerate}
\def\labelenumi{\arabic{enumi}.}
\tightlist
\item
  複素極表示は、位相角を含む極表示を直交二成分へ写す標準写像である。
\item
  その共役積は二つの直交実成分の二乗和であり、平方量読出しの一形式として読める。
\item
  実成分を二本ずつ束ねれば、対等な多次元二乗和
  \(\sum_j Z_j\bar Z_j=\sum_n x_n^2=R^2\) に一般化できる。
\item
  複素平面に表示されない不可視軸も、同符号の対等な軸として加わるだけであり、半径
  \(R\) は一意である。
\item
  共役積 \(X\bar X\) と非共役の二乗 \(X^2\)
  は、実座標では一致するが虚座標でのみ分岐する（\(X_h\bar X_h=z_h^2\)
  に対し \(X_h^2=-z_h^2\)）。負号はこの分岐からのみ生じる。
\item
  共役ノルム背骨に虚座標二乗を加えた零二乗和形式
  \(\sum_j Z_j\bar Z_j+\sum_k (i z_k)^2=0\)
  を展開すると、実座標の二乗和が虚座標ラベルの二乗和に等しい、という代数関係
  \(\sum_n x_n^2=\sum_k z_k^2\)
  が現れる。非自明であるには少なくとも一つの虚数値座標を要する。
\item
  これらはすべて、共役複素数と二乗和の代数的配置に関する幾何学的観察であり、物理学の主張を含まない。
\end{enumerate}

\begin{center}\rule{0.5\linewidth}{0.5pt}\end{center}

\section{9. 結論}\label{ux7d50ux8ad6}

本稿では、平方量読出し・奇数倍音位相構造・相対位相コントラストの系列に続く第四篇として、共役複素ノルムと平方量読出しの接続を観察し、それを対等な多次元二乗和へ一般化した。

複素極表示 \(Z=\rho e^{i\theta}\) に
\(x=\rho\cos\theta,\ y=\rho\sin\theta\) を置けば
\(Z\bar Z=x^2+y^2=\rho^2\)
である（標準的事実）。本稿の中心は、この既知のノルムを平方量読出し系列に置き、さらに実成分を二本ずつ束ねて対等な多次元二乗和

\[
\sum_j Z_j\bar Z_j=\sum_n x_n^2=R^2
\]

に一般化した点にある。ここに特別な軸はなく、不可視軸も同符号の対等な軸として加わり、\(R\)
は一意である。

さらに、この共役ノルム背骨に虚数値座標の二乗を加えた零二乗和形式

\[
\sum_j Z_j\bar Z_j+\sum_k (i z_k)^2=0
\]

を観察した。共役積 \(X\bar X\) と非共役の二乗 \(X^2\)
は実座標では一致し虚座標でのみ分岐する（\(X_h\bar X_h=z_h^2\) vs
\(X_h^2=-z_h^2\)）ため、展開すると純虚数値座標の \(i^2=-1\) を通じて

\[
\sum_{n} x_n^2=\sum_{k} z_k^2
\]

という代数関係が現れる。ここで現れる負の二乗寄与は \(i^2\)
の帰結であって、手で置いた計量の署名ではなく、\(R\)
の二重定義も生じない。共役ノルム型 \(Z\bar Z\) と虚座標の非共役二乗
\(X_h^2\) は同一の演算ではなく、本稿はこの区別を保つ。

以上は、共役正定値ノルム \(\sum_j Z_j\bar Z_j\)
と、虚数値座標に対する非共役二乗項 \((i z_k)^2\)
を、同一の零二乗和形式の中でどのように接続できるかを、平方量読出し・対等な多次元二乗和との整合とともに記録した、初等代数と幾何的解釈に限った観察である。本稿は物理学（時空・相対論・量子論・計量署名）とは接続せず、座標ラベルを物理的次元と同定しない。

\begin{center}\rule{0.5\linewidth}{0.5pt}\end{center}

\section{参考文献}\label{ux53c2ux8003ux6587ux732e}

\subsection{自己参照}\label{ux81eaux5df1ux53c2ux7167}

{[}1{]} 木原範昭,
「平方量読出しによる線形単体化と曲率補正候補の次元別整理」, v0.1-r1,
Zenodo DOI: 10.5281/zenodo.20785540, 2026.

{[}2{]} 木原範昭,
「半波長位相区間における一定振幅奇数倍音和の孤立ピーク波とその局在性に関する観察」,
v0.6, Zenodo DOI: 10.5281/zenodo.21073985, 2026.

{[}3{]} 木原範昭,
「半波長奇数倍音孤立ピーク波の二コピー共通相対位相重ね合わせにおける波形不変性とコントラスト則の観察」,
v0.1, Zenodo DOI: 10.5281/zenodo.20923462, 2026.

\subsection{関連する標準事項}\label{ux95a2ux9023ux3059ux308bux6a19ux6e96ux4e8bux9805}

{[}4{]} 高木貞治『解析概論』改訂第3版、岩波書店、1961年.

{[}5{]} 複素数の極形式、Euler
公式、複素ノルム、および標準的な二次形式に関する一般的教科書事項。

\begin{center}\rule{0.5\linewidth}{0.5pt}\end{center}

\section{付録A:
式の最小連鎖（修正版）}\label{ux4ed8ux9332a-ux5f0fux306eux6700ux5c0fux9023ux9396ux4feeux6b63ux7248}

\begin{enumerate}
\def\labelenumi{\arabic{enumi}.}
\tightlist
\item
  複素平面で \(Z=x+iy\) と置く（\(i\) の役割①：束ね）。
\item
  共役積により \(Z\bar Z=x^2+y^2=\rho^2\)。
\item
  実成分を二本ずつ束ね、対等に加える：\(\displaystyle\sum_j Z_j\bar Z_j=\sum_n x_n^2=R^2\)。不可視軸も同符号の対等な軸として加わるだけ。
\item
  共役積と二乗の違いを確認：実座標では
  \(x\bar x=x^2\)、虚座標でのみ分岐し \(X_h\bar X_h=z_h^2\) に対し
  \(X_h^2=(i z_h)^2=-z_h^2\)（\(i\) の役割②）。負号はここからのみ。
\item
  共役ノルム背骨に虚座標二乗を加えた零二乗和形式：\(\displaystyle\sum_j Z_j\bar Z_j+\sum_k (i z_k)^2=0\)。非自明性が虚数値座標を要求する。
\item
  展開・移項：\(\displaystyle\sum_j Z_j\bar Z_j=\sum_k z_k^2\)、すなわち
  \(\displaystyle\sum_{n} x_n^2=\sum_{k} z_k^2\)。負号は \(i^2\)
  の帰結。\(Z\bar Z\) 型と \(X^2\) 型を混同せず、\(R\)
  の二重定義も生じない。
\end{enumerate}

\begin{center}\rule{0.5\linewidth}{0.5pt}\end{center}

\section{付録B:
次論文への接続候補（本稿の範囲外）}\label{ux4ed8ux9332b-ux6b21ux8ad6ux6587ux3078ux306eux63a5ux7d9aux5019ux88dcux672cux7a3fux306eux7bc4ux56f2ux5916}

本稿は、代数的配置に関する観察に限定した。次論文で検討しうる候補（いずれも本稿では扱わない）：

\begin{enumerate}
\def\labelenumi{\arabic{enumi}.}
\tightlist
\item
  虚数値座標を導入する幾何学的動機の整理（どの軸を虚に置くかの選択規準）。
\item
  零二乗和 \(\sum_n X_n^2=0\) の解集合（複素二次超曲面）の幾何。
\item
  実部分空間への射影と、不可視軸の役割の一般化。
\item
  中心投影シリーズの投影半径との比較（比較のみ。同定はしない）。
\item
  物理的解釈へ進む場合に必要となる追加仮定の明示（本稿はここへ進まない）。
\end{enumerate}

\begin{center}\rule{0.5\linewidth}{0.5pt}\end{center}

\section{付録C:
共役対閉包と非復元性の入口}\label{ux4ed8ux9332c-ux5171ux5f79ux5bfeux9589ux5305ux3068ux975eux5fa9ux5143ux6027ux306eux5165ux53e3}

本文 §5 では、共役ノルム型平方量 \(\sum_j Z_j\bar Z_j\)
に虚数値座標の非共役二乗 \((i z_k)^2=-z_k^2\) を加えて、非自明な零二乗和
\(\sum_j Z_j\bar Z_j+\sum_k (i z_k)^2=0\)
を得た。本付録は、この構成を逆に読んだときに現れる、ごく単純な観察を記録する。\textbf{不確定性原理を導出するものではない。}

\subsection{C.1
実共役対だけなら、零二乗和は自明に閉じる}\label{c.1-ux5b9fux5171ux5f79ux5bfeux3060ux3051ux306aux3089ux96f6ux4e8cux4e57ux548cux306fux81eaux660eux306bux9589ux3058ux308b}

成分がすべて実数値で、固定された組により
\(Z_j=q_{2j-1}+i\,q_{2j}\)（\(q_n\in\mathbb{R}\)）と余りなく束ねられるなら、

\[
\sum_j Z_j\bar Z_j=\sum_n q_n^2\ge0
\]

は正定値であり、\(\sum_n q_n^2=0\) は全零解 \(q_1=\cdots=q_{2m}=0\)
しか許さない。すなわち、この範囲では零二乗和は\textbf{自明に閉じる}（非自明な解を持たない）。

\subsection{C.2
非自明な零二乗和は非実成分を要する}\label{c.2-ux975eux81eaux660eux306aux96f6ux4e8cux4e57ux548cux306fux975eux5b9fux6210ux5206ux3092ux8981ux3059ux308b}

したがって、非自明な零二乗和を得るには、\((i z_k)^2=-z_k^2\)
のような負寄与、すなわち\textbf{少なくとも一つの非実成分}（または固定共役対構造の破れ）が必要である（§5.3）。言い換えれば、実共役対だけでは正定値で\textbf{閉じすぎて}、非自明なゼロを作れない。

\subsection{C.3
共役ノルム読出しは位相を潰す（多対一）}\label{c.3-ux5171ux5f79ux30ceux30ebux30e0ux8aadux51faux3057ux306fux4f4dux76f8ux3092ux6f70ux3059ux591aux5bfeux4e00}

一方、共役ノルム読出しそのものは、\textbf{閉じているか否かに関わらず}位相を落とす。\(Z=\rho e^{i\theta}\)
に対し \(Z\bar Z=\rho^2\) であり、\(\theta\) は消える。同じ \(\rho\)
を持つすべての \(Z=\rho e^{i\theta}\) が同じ \(\rho^2\) を返すので、写像

\[
Z\longmapsto Z\bar Z=\rho^2
\]

は\textbf{多対一}である。したがって、二乗量 \(\rho^2\) から元の位相表示
\(Z\) を一意に復元することはできない。

\subsection{C.4
非復元性の入口}\label{c.4-ux975eux5fa9ux5143ux6027ux306eux5165ux53e3}

C.1--C.3
を合わせると、二乗量読出しには\textbf{二つの相補的な側面}がある。

\begin{itemize}
\tightlist
\item
  \begin{enumerate}
  \def\labelenumi{(\roman{enumi})}
  \tightlist
  \item
    \textbf{零の自明性}：実共役対に閉じる限り、零二乗和は正定値で自明であり、非自明化には非実成分を要する（C.1,
    C.2）。
  \end{enumerate}
\item
  \begin{enumerate}
  \def\labelenumi{(\roman{enumi})}
  \setcounter{enumi}{1}
  \tightlist
  \item
    \textbf{読出しの多対一性}：共役ノルム読出しは、閉包の有無に関わらず、常に位相方向を同一視する（C.3）。
  \end{enumerate}
\end{itemize}

本稿はこの (ii)
の多対一性を、\textbf{非復元性の入口}と呼ぶ。これは「隠れた真の位相があるのに観測できない」という隠れた変数の主張ではなく、\textbf{読出し写像がそもそも位相方向を区別しない}という縮退の観察である。位相消去は閉包の破れによって生じるのではなく、共役ノルム読出しに内在する点に注意する。

\subsection{C.5
位置づけと非主張}\label{c.5-ux4f4dux7f6eux3065ux3051ux3068ux975eux4e3bux5f35}

本付録は、量子力学の不確定性原理を導出せず、それを置換せず、非可換演算子形式の代替でもない。実
Hilbert 空間に複素構造 \(J^2=-1\)
を入れると複素量子論が現れること、概複素構造が偶数次元性と結びつくこと、Born
則 \(\psi\mapsto|\psi|^2\)
で位相情報が落ちることは、いずれも既知の標準事項である。本付録が付け加えるのは、これらの\textbf{前段}にある「成分が固定された共役対として閉じるか」という観点を、二乗量読出しの決定性と非復元性を見分ける\textbf{入口}として明示することだけである。座標ラベルを物理的次元と同定せず、\(z_k\)
を時間・曲率半径と同一視しない立場（§0, §6）は本付録でも保持する。
\end{document}
